\cotefile{94}
\scanpage{2}
\uline{Construction des} \uline{Modules cohérents}

\deleted{et de schémas \unsure{fibrés|finis} sur d'autres par descente} par leur foncteur représentatif

\noindent\hfill\rule{1.3cm}{0.5pt}\hfill\mbox{}

Application à la proreprésentation d'un foncteur $G\deleted{\unsure{\to|-}}/S$ le long d'une section, à la définition des invariants $\Omega^{1}\,\unsure{\mathfrak{g}|g|(y)}$, \ldots.

% archive sheet number: 1
\scanpage{3}
\margin{$\big/\!\big/$}\uline{Construction de \unsure{faisceau|faisceaux} de modules}

$A$ anneau \deleted{\unsure{noeth.}}
\[
F:\operatorname{Mod}(A)\longrightarrow(\mathit{Ens})\qquad\text{foncteur}
\]
Pour toute algèbre $B$ sur $A$, on en conclut un foncteur
\diagram{\begin{tikzcd}[ampersand replacement=\&,row sep=small]
F_B:\operatorname{Mod}(B) \arrow[rr] \arrow[dr,"\mathrm{restr.}"'] \& \&(\mathit{Ens})\\
\&\operatorname{Mod}(A)\arrow[ur]\&
\end{tikzcd}}
Si $F$ repr. par $\deleted{\unsure{M|N}}P$, $F_B$ l'est par $\deleted{\unsure{M|N}}P_B$

Si $f\in A$, on pose $F_f$ au lieu de $F_{A_f}$.

\uline{Lemme 1} \deleted{Supposons $F$ exact à gauche, et soit} \inserted{Soit $B$ une $A$-algèbre fid. plate.} Alors
\[
\begin{aligned}
&F\deleted{{}_{\unsure{f}}}\ \text{représentable (resp. par \inserted{$A$-}mod de type fini, de prés. finie)}\\
&\Updownarrow\\
&F\ \text{exact à gauche, et } F_B\ \text{représentable \deleted{par} (resp. par \inserted{$B$-}module}\\
&\hspace{12em}\text{de type fini, de prés finie).}
\end{aligned}
\]

\deleted{\uline{Corollaire 2}} \inserted{Lemme 2} Soit $f\in A$, supposons $A$ noeth. Alors $F$ repr. \deleted{\unsure{ssi}} par \deleted{un} $A$-module de t.f. ssi il satisfait :

(i) $F$ exact à gauche

(ii) $F(\overset{\deleted{\wedge}}{M})\xrightarrow{\sim}\varprojlim F(M/f^nM)$ \quad si $M$ un $\widehat A$-module \deleted{$(=\widehat M=\varprojlim M/f^nM$} de t.f. \fbox{$\widehat A=\varprojlim A/f^nA$}

(iii) $F$ commute aux $\varinjlim$ filtrantes

(iv) $F_f$ repr. \deleted{\unsure{par $A$}}

(v) $\forall n,\quad F_{A/f^nA}$ représentable.

Nécessité claire, (iii) exprimant le fait que le module $P$ qui représente $F$ est de t.f. (donc de prés. finie).

% archive sheet number: 2
\scanpage{4}
Suffisance. Grâce à \textcircled{\scriptsize i} et Lemme 1 on est ramené à voir que $F_{\widehat A}$ est représentable, car posant $B=\widehat A\times A_f$, $F_B$ sera alors représentable en vertu de \textcircled{\scriptsize iv}, donc $F$ l'est ; le $P$ le représentant sera alors de t.f. grâce à \textcircled{\scriptsize iii}.

Pour tt $n$, $F_{\underbrace{A/f^{n+1}A}_{A_n}}$ est représentable, par hyp. \inserted{\textcircled{\scriptsize{v}}} par un module $P_n$ sur $A_n$, \deleted{Soit $P$} de type fini grâce à (iii). Soit $\widehat P=\varprojlim P_n$, qui est de t.f. sur $\widehat A$ (N.B. les $P_n$ ``se recollent''). \deleted{Pour \unsure{tout module}} Utilisant \textcircled{\scriptsize ii}, on trouve que $\widehat P$ représente $F_{\widehat A}$ sur la catégorie des $A$-modules de t.f., et utilisant \textcircled{\scriptsize iii}, on trouve que $\widehat P$ représente $F_{\widehat A}$ partout.

\uline{Proposition} 3 \quad $A$ noeth. Pour que $F$ soit repr. par un module \deleted{de} t.f. il f. et s. que $F$ satisfasse aux conditions suivantes :

(i) $F$ exact à g.

(ii) $F$ commute aux $\varinjlim$ filtrantes.

% A bracket encloses the last physical condition; its arrowhead falls between (ii) and the condition on quotients.
\margin{\tikz{\draw[->] (0,0) .. controls (-.5,.6) and (-.45,1.05) .. (.12,1.16);}}
\inserted{(\deleted{iv}iii) Pour tout $f\in A$, et tout module de type fini $M$ sur $\widehat A$\inserted{$=\varprojlim A/f^{n+1}A$}, \unsure{on}
\[
F(\deleted{\widehat M}M)\xrightarrow{\sim}\varprojlim F(M/f^{n+1}M)
\]
\deleted{où $\widehat M=\varprojlim M/f^{n+1}M$.}}

(\deleted{iii}iv) Pour tout quotient \inserted{non nul} $B$ de $A$, \deleted{et tout $f\in B$,} \inserted{il existe un} $f\in B$, non nilpotent, tel que $F_{B_f}$ soit représentable.

% archive sheet number: 3
\scanpage{5}
Nécessité claire. Pour la suffisance, on note que les conditions sont stables par passage à une algèbre quotient $B$. Par récurrence noeth, on peut donc supposer que pour tt tel quotient $B=A/J$, avec $J\ne0$, $F_B$ est représentable (par module de t.f.). On peut supposer aussi $A\ne0$. Il existe par \textcircled{\scriptsize iv} un $f\in A$ non nilpotent, tels que $F_{\deleted{\unsure{f}}A_f}$ soit repr. D'autre part, par l'hyp. de récurrence noeth, comme pour tt $n$, on a $f^{n+1}\ne0$, on voit que $F_{A/f^{n+1}A}$ est repr. On peut donc appliquer le lemme 2, qui donne la conclusion.

\noindent\rule{\linewidth}{0.4pt}\par\nobreak\noindent\rule{\linewidth}{0.4pt}

En pratique, la condition (iv) est parfois difficile à vérifier. On va donner d'autres conditions.

\uline{Proposition 4} \quad $A$ \uline{local} \uline{noeth}. Pour que $F$ soit repr. par module de t.f., il f. et s. que $F$ satisfasse les conditions

(i) $F$ exact à gauche

(ii) $F$ commute aux $\varinjlim$

(iii) Pour tt $\widehat A$-module de t.f. $M$,
\[
F(\widehat M)\xrightarrow{\sim}\varprojlim_{n}F(M/{\unsure{\mathfrak m|m}}^{n+1}M),\ \text{et}
\]

(iv) \deleted{\unsure{de}} $F(\widehat M)$ est de t.f. sur $\widehat A$ (il suffit $=$ si $M$ est \uline{fini} sur $A$).

Nécessité claire. Pour suffisance, grâce au Lemme 1 et condition (i), on peut supposer $A$ complet.

% archive sheet number: 4
\scanpage{6}
En vertu de (i), sa restriction aux modules \deleted{\unsure{artiniens}} \inserted{\uline{finis} sur $A_n=A/{\unsure{\mathfrak m|m}}^{n+1}$} est proreprésentable, \deleted{par $(P_n)$ ($P_n$ \unsure{module} sur $A_n$,} \deleted{\illegible{1}} et en vertu de la condition (iv), \unsure{testant|testée} avec l'injectif indécomposable $I_n$ de $A_n$, on trouve qu'il est $=$ représentable, \deleted{\unsure{d'où ses|\illegible{1} par} $P_n$,} par un module de t.f. $P_n$ sur $A_n$. Les $P_n$ se recollent, soit $P=\varprojlim P_n$, qui est un $A$-module de t.f. On conclut de (iii) que $P$ représente $F$ sur les modules de t.f., et en vertu de (ii) sur tous les modules, cqfd.

\deleted{\uline{Proposition}} \inserted{Lemme 5} \quad $A$ noeth. \deleted{Conditions équivalentes}

% A long arrow carries the bracketed statement from here to the place after « Alors ».
\deleted{1)} \margin{\tikz[baseline]{\draw[->] (0,1) .. controls (-.7,.5) and (-.7,-.5) .. (0,-1);}}

\deleted{2) \unsure{On}} \inserted{Supposons qu'on a} les conditions

(i) $F$ exact \deleted{\unsure{est}} à g

(ii) $F$ commute aux $\varinjlim$ filtrantes

(iii) $\forall\,p\in\operatorname{Spec} A$, et tt $M$ de t.f. $\widehat A_p$, on a
\[
F(M)\xrightarrow{\sim}\varprojlim F(M/\deleted{\unsure{p}} \mathfrak m_p^{n+1}M)
\]

(iv) $\forall$ module de \deleted{type fini} \inserted{longueur finie} $M$ sur $A_p$, $F(M)$ est \unsure{.}\ un $A_p$-module de t.f.

(v) $\forall$ idéal premier $p$, \deleted{$\unsure{\forall}$} $\exists\, f\in\unsure{A/p|A-p|A/\mathfrak p}$ tels que $F_{(A/p)_f}$ soit représentable

Alors \inserted{$F$ proreprésentable par un syst. projectif strict \inserted{$(P_i)$} de modules de t.f., tel que pour tout $p\in\operatorname{Spec} A$, le syst. des $(P_i)_p$ soit stationnaire (i.e. $F_p$ est représentable par un module de t.f. sur $A_p$).}

\uline{Dém}. Pour tout $p\in S=\operatorname{Spec} A$, la prop. 4 implique que $F_p$ est représentable par $P_p$ de t.f. sur $A_p$. (grâce à (i) à (iv)). Les $P_p$ sont stables par

% archive sheet number: 5
\scanpage{7}
\unsure{générisation|généralisation}. Pour achever de prouver la prop, il suffit de prouver (compte tenu de \textcircled{\scriptsize i} \inserted{et \textcircled{\scriptsize ii}}) que pour tout $\alpha\in F(M)$, $M$ de t.f. sur $A$, il existe un ``module de définition $N$'' i.e un plus petit sous-module $N$ de $M$ tel que $\alpha\in\operatorname{Im}(F(N)\to F(M))$. Or on a pour tt $p\in S$ un $N_p\subset M_p$ $(=\operatorname{Im}(\alpha_p:P_p\unsure{-|\to} M_p))$. \deleted{Soit tt} Il suffit de voir que le système des $N_p$ provient d'un sous-module $N$ de $M$ (on voit alors aisément qu'il conviendra). Soit $U\subset S$ le plus grand ouvert sur lequel ça marche, \inserted{par un $\underline{\widetilde N}_U\subset\widetilde M|U$} Si on avait $U\ne S$, soit $x\in S-U$ \deleted{\unsure{un pt maxi|complètement}} un pt maximal de $S-U$.

\deleted{Il existe un voisinage \unsure{ouvert}} \deleted{Il existe un} Soit $S'=\operatorname{Spec}\mathcal O_{S,x}$, $U'=U\cap S'$, \deleted{donc} $N_{S'}=\widetilde{N_x}$ (NB. \tikz[baseline=(a.base)]{\node[inner sep=1pt](a){$x=p$};\draw (a.center) ellipse (.48cm and .23cm);} !), alors $\widetilde N_x|U'=\underline N_U|U'$, donc $\exists$ voisinage ouvert $V$ de $x$, \deleted{tel que} et un $\underline{N}'\subset\widetilde M|V$, tel que $\underline N'_x=N_x\subset M_x$ et $\underline N'|V\cap U=\underline N_U|V$, i.e.
\[
\left\{\begin{aligned}
\underline N'_x&=N_x\subset M_x\\
\underline N'_s&=N_s\subset M_s\qquad\text{si }\deleted{\unsure{x}}s\in U\cap V.
\end{aligned}\right.
\]
\margin{$\alpha_V\in F(N')$ donc $N_s\subset\underline N'_s$ pour tout $s\in V$}
\margin{$\longrightarrow$ Donc on est ramené au cas $\underline N'=\unsure{\widetilde M|M}$ \deleted{\unsure{quitte à}}}
Prenant $V$\inserted{$=\operatorname{Spec} A_f$} assez petit, \deleted{la condition (v) implique} \inserted{on peut supposer que} \deleted{\unsure{et}} la condition (v) implique qu'on peut supposer que $F_{(A/p)_f}$ est représentable\unsure{?|.}\ \deleted{Il s'ensuit que} \deleted{et} On voit tout de suite que ceci implique que (quitte à restreindre encore $V$) l'image de $N_s$ dans $M_s$ \deleted{est \uline{égale} \unsure{déjà}} \inserted{est tout} modulo $p$, donc est tout pour $s\in V(p)$ par Nakayama. \unsure{Ceci gagne|Ceci signifie|C'est gagné|Cela suffit}.

% archive sheet number: 6
\scanpage{8}
\uline{Remarque} \uline{5 bis} \quad Avec les notations du lemme 5, et en présence des conditions (i) à (iv), la condition (v) équivaut à la condition suivante, qui se prête mieux à une vérification directe :

(v bis) La fonction $p\longmapsto\operatorname{rang}_{k(p)}F(k(p))$ \deleted{\unsure{sur|sera}} sur $S$ est constructible. \inserted{Il faut cependant supposer que $A$ satisfait à la condition : pour tout quotient intègre $B$ de $A$, $\unsure{\operatorname{Nor}|\operatorname{Nr}}(B)$ contient un ouvert non vide.}

Le fait que $(\mathrm v)\Longrightarrow(\mathrm v\ \mathrm{bis})$ est trivial.

En sens inverse, on utilise le fait que $F$ définit un ``syst. local'' de modules $\deleted{\unsure{M|Q}}P_p$. Dans le cas $A$ intègre, \deleted{l'hyp.} si le rang des $\deleted{M}P_p\otimes k(p)$ est constant et \unsure{égale|égal} à $n$ pour $p\in U$ (ouvert non vide), \deleted{et \unsure{x}} on en conclut que les $\deleted{M}P_p$ $(p\in U)$ sont alors \uline{libres} de rang $n$. D'ailleurs, \deleted{\unsure{de par}} comme $F$ commute aux $\varinjlim$, on voit qu'on peut supposer que les $\deleted{M}P_p$ sont tous $\subset N$, avec $N$ libre de rang $\unsure{n|N}$. Utilisant le fait que \deleted{\unsure{alors le}} $F(\deleted{\unsure{M}}N)$ est de type fini, on trouve aisément que pour tous les $p\in U$ de codim. 1, \uline{sauf un nombre fini}, on a $P_p=N_p$. Quitte à restreindre $U$, on peut supposer que ceci est vrai pour tous les $p$ de codim 1 de $U$. Mais alors, si on prend $U$ assez petit pour être normal, on conclut que $P_p=N_p$ pour \uline{tt} $p\in U$, \deleted{donc} d'où facilement que $N$ représente $F|U$.

% archive sheet number: 7
\scanpage{9}
\uline{Lemme 6} \quad Soient $S$ schéma \deleted{\unsure{de}} \inserted{$\mathrm{loc}^{\unsure{t|.}}$} noeth, $F$ \deleted{\illegible{1}} coh. sur $S$.

a) Soit $J$ le Nilradical de $\mathcal O_S$, \inserted{$S_0=S_{\mathrm{réd}}$.} Il existe un ouvert dense $U$ de $S$ tel que \deleted{\illegible{3} \unsure{et|(2)} $\unsure{\operatorname{gr}|\operatorname{Supp}}_{\!J}(F)$} $\operatorname{gr}_J(\mathcal O_S)$ \inserted{et $\operatorname{gr}_J(F)$} loc. libres sur $U_0$

b) Soit $U$ comme dessus. Alors pour tout $s\in U$, et \deleted{toute} suite $\mathcal O_{S,s}$-régulière $f_1\ldots f_n$ \deleted{de} $\unsure{\mathfrak m_s|M_s}$, \deleted{posant $M=\mathcal O_{S,s}/(f_1,\ldots,f_n)$,} on a
\[
\operatorname{Tor}^{\mathcal O_{S,s}}_i\bigl(\deleted{\unsure{\mathcal O_{S,s}}}F_s,\mathcal O_{S,s}/(f_1,\ldots,f_n)\bigr)=0\qquad\text{pour }i>0.
\]

\uline{Dém} L'assertion a) est triviale. Pour b), on est ramené au cas où $S=U$, \deleted{et $F$ provient d'un faisceau loc libre sur $S_0$.} \inserted{et où $S=\operatorname{Spec}\mathcal O_{S,s}$,} Pour calculer les $\operatorname{Tor}_i$ on utilise le complexe de chaines de Koszul $K_{\bullet}(\underrightarrow f)$, qui est une résolution de $A/(f_1\ldots,f_n)$ $(A=\mathcal O_{S,s})$. On est ramené à prouver que les images $f_{10}\ldots f_{n0}$ des $f_i$ de $A_0$ forment \deleted{\unsure{encore|aussi}} \inserted{aussi} une suite $M$-régulière. En fait, on prouve en $=$ temps que, posant $N=M/(f_1\ldots f_n)M$, on a $\operatorname{gr}_J(N)$ \deleted{loc. libre} libre. C'est bien standard si $n=1$, et on continue par récurrence.

\uline{Corollaire} 7 \quad $S$ schéma noethérien, $P$ et $\Omega$ faisceaux cohérents. Alors $\exists$ ouvert dense \inserted{$U$} de $S$ tel que pour tt $s\in U$ et tt suite $\underrightarrow f$ $\mathcal O_{S,s}$-régulière $f_1\ldots f_n$,

% archive sheet number: 8
\scanpage{10}
l'hom. naturel
\[
\underline{\operatorname{Hom}}(P,\Omega)_s\otimes\mathcal O_{S,s}/(f_1\ldots f_n)\longrightarrow
\underline{\operatorname{Hom}}_{\mathcal O_{S,s}}\bigl(P_s,\Omega_s/(f_1\ldots f_{\unsure{s|n}})\Omega_{\unsure{s}}\bigr)
\]
soit un isom.

\uline{Dém} On peut supposer $S$ affine \inserted{anneau $A$}, alors on a \deleted{résolution} présentation
\[
L_1\mathbin{-}L_0\longrightarrow P
\]
avec $L_1,L_0$ \unsure{libre|libres} de t.f., donc $\underline{\operatorname{Hom}}(P,Q)$ est le $H^0$ du complexe
\[
K^{\cdot}:\qquad 0-\overset{\text{degré }0}{\check L_0\otimes P}-\overset{\text{degré }1}{\check L_1\otimes P}-0
\]
tandis que $\operatorname{Hom}(P,Q\otimes M)$ est le $H^0$ du complexe $K^{\cdot}\otimes M$. Un critère pour que l'hom naturel
\[
H^0(K^{\cdot})\otimes M^{\unsure{\cdot}}-\ H^0(K^{\cdot}\otimes M)
\]
soit bijectif est qu'on ait
\[
\operatorname{Tor}_1(M,H^1(K^{\cdot}))=\operatorname{Tor}_1(M,Z^1(K^{\cdot}))=0
\]
Or $H^1(K)=R$ et $Z^1(K)=S$ sont des modules de t.f. sur $A$, auxquels on peut appliquer le \deleted{corollaire} lemme 6, b), en faisant $M=\mathcal O_{S,s}/\unsure{(f_1\ldots f_n)|f_1\ldots f_n}$.

On gagne.

% archive sheet number: 9
\scanpage{11}
$\Omega$ module de différentielles pour un C.M. $\Longleftrightarrow$ les $\Omega/(f_1\ldots f_n)\Omega$ sont les injectifs indécomposables

\deleted{\uline{Lemme \unsure{1}} Soit $X$ un préschéma régulier, \unsure{un} \illegible{1} \unsure{fibré}} \deleted{\unsure{fini}}

\uline{Lemme 8} Soit $X$ un préschéma de Cohen-Macaulay ; \margin{\tikz{\draw (0,0)--(0,1.8); \draw (.08,0)--(.08,1.8);}} $\Omega$ un faisceau fondamental sur $X$, alors pour tte suite $(f_1\ldots f_n)$ de sections de $\mathcal O_X$ qui est $\mathcal O_X$-régulière, $\Omega/(f_1,\ldots f_n)\Omega$ est un faisceau fondamental sur le \deleted{\unsure{pré}}schéma \deleted{\unsure{noethérien|noeth.}} \inserted{CM} \deleted{$\unsure{H^0(|H^2(}$} $V(f_1,\ldots f_n)$.

\uline{Corollaire} 9. Soit $x\in X$ et $f_1\ldots f_n\in\mathcal O_{X,x}$ un syst. de paramètres de $\mathcal O_{X,x}$, alors \margin{\tikz{\draw (0,0)--(0,1.7);}} $\Omega/(f_1\ldots f_n)\Omega$ est \deleted{\unsure{un}} l'injectif indécomposable dans la catégorie des modules sur $\mathcal O_{X,x}/(f_1\ldots f_n)$.

\uline{Corollaire 10} Soit $X$ un préschéma loc. noeth \margin{\tikz{\draw (0,0)--(0,1.5);}} qui est \uline{localemt immergeable dans} un \uline{régulier}.
% slanted
\margin{il suffit \deleted{\unsure{qu'il|que}} qu'il soit excellent, ou seulement qu'il satisfasse les conditions de EGA \uline{IV} 6.11.8.}
Alors il existe un ouvert dense $U$ \deleted{de} $X$, et un Module coh. $\Omega$ sur $X$, ayant les propriétés suivantes

a) Les anneaux locaux de $U$ sont C.M.

b) Pour tout $x\in U$, et tt syst. de paramètres $f_1\ldots f_n$ de $\mathcal O_{X,x}$, le module $\Omega/(f_1\ldots f_n)\Omega$ sur $\mathcal O_{X,x}/(f_1,\ldots f_n)$ est

% archive sheet number: 10
\scanpage{12}
\margin{-}\margin{-}CHAPITRE I

ARITHM\handwritten{\'}{E}TIQUE DES CORPS VALU\handwritten{\'}{E}S .

\S~1 . \deleted{\unsure{C}}Valeurs absolues et valuations .

Rappelons qu'étant donné un corps $K$ , on appelle \uline{valeur absolue} sur $K$ une application $f$ de $K$ dans $R_{+}$ telle que $f(xy)=f(x)f(y)$ , $f(x+y)\leqslant f(x)+f(y)$ ,et que $f(x)=0$ entraîne $x=0$ ; la distance $f(x-y)$ définit alors sur $K$ une topologie compatible avec sa structure de corps .

Soit maintenant $v$ une valuation de $K$ dont le groupe des ordres est archimédien , c'est-à-dire isomorphe à un sous-groupe additif de $R$ . La fonction numérique $e^{-v(x)}$ est une valeur absolue sur $K$ ; on dit qu'une telle valeur absolue est \uline{valuative} .

PROPOSITION 1 .- \uline{Pour qu'une valeur absolue $f$ d'un corps} $K$ \uline{soit valuative} , \uline{il faut et il suffit que l'une ou l'autre des conditions suivantes soit vérifiée} :

a) $f(x+y)\leqslant \operatorname{Max}(f(x),f(y))$ ;

b) $f(n.1)\leqslant 1$ \uline{pour tout entier naturel $n$} .

C'est évident pour a) . a) entraîne b) car $f(1)=1$ . Réciproquement la formule du binôme montre que $f((a+b)^n)\leqslant (n+1)\operatorname{Max}(f(a)^n,f(b)^n)$ , d'où a) en faisant tendre $n$ vers l'infini .

COROLLAIRE \deleted{1} .- \uline{Toute valeur absolue d'un corps de caractéristique} $p\neq0$ \uline{est valuative} .

En effet , tout élément $n.1$ est une racine de l'unité .

On dit (\uline{Top.gén}.,chap.IX,\S~3\deleted{\unsure{1|4}},n\textsuperscript{o}2) que deux valeurs absolues $f,g$ de $K$ sont \uline{équivalentes} s'il existe un nombre réel $s$ tel que $f(x)=g(x)^s$ .

PROPOSITION 2 .- \uline{Toute valeur absolue $f$ de} $Q$ \uline{est équivalente à une valeur absolue} $p$-\uline{adique , ou à la valeur absolue ordinaire} .

Ceci a déjà été vu dans le cas où $f$ est valuative (chap. des valuations) .

Sinon , soit $n$ un entier tel que $f(n)>1$ . Soient $a$ et $b$ des entiers \inserted{\handwritten{$>1$}} ; écrivons $\deleted{\unsure{a^{b}=a^{0}|a^{b}=a_{0}}}\;a^{h}=a_{0}+a_{1}b+\ldots+a_{q}b^{q}$ , où $0\leqslant a_{i}\leqslant b-1$ . On en déduit

% archive sheet number: 11
\scanpage{13}
l'injectif indécomposable

Ces conditions impliquent la suivante

c) Pour tout hom $u\colon R\mathbin{-}R'$ surjectif de modules de type fini sur un $\mathcal O_{\deleted{\unsure{M|A}}X,x}$ $(x\in U)$, tel que $\unsure{\operatorname{Hom}|\operatorname{Rhom}}(R',\Omega_x/(f_1\ldots f_n)\Omega_x)\to\operatorname{Hom}(R,\Omega_x/(f_1\ldots f_n)\Omega_x)$ soit un isomorphisme, \inserted{pour tt syst. de paramètres $f_1\ldots f_n$ de $\mathcal O_{X,x}$,} \deleted{$\unsure{e|c}$}$u$ est \unsure{lui $\overline{\mathrm m}$|lui-m} un isomorphisme

% Several diagonals cancel the Corollaire 11 draft below. Horizontal strokes and local rewrites are distinguished inside the cancellation.
\deleted{\uline{Corollaire 11} \deleted{Supposons \unsure{que}} \inserted{Supposons que} $U,\Omega$ soient comme dans Cor. 10. Soit $x\in U$, et soit $u\colon R\mathbin{-}R'$ un hom \inserted{\uline{surjectif}} de modules de type fini sur $\mathcal O_{X,x}$, \deleted{qui soit} tel que $R$ et $R'$ satisfassent la condition suivante pour les foncteurs $\Phi_R(M)=\deleted{\operatorname{Hom}(R,\unsure{D\otimes M|D M|\Omega\otimes M})}\inserted{\operatorname{Hom}(R,M)}$ et $\Phi_{R'}(M)=\deleted{\operatorname{Hom}(R',\unsure{D\otimes M|D M|\Omega\otimes M})}\operatorname{Hom}(R',M)$ :}

% slanted
\margin{à expliciter d'une autre façon \tikz{\draw (0,0)--(0,1.1);}}

\deleted{a)}
\[
\deleted{\deleted{\unsure{\Phi_R|\Phi}}\Phi(\Omega)\otimes A/(\deleted{\unsure{f}}f_1,\ldots f_n)\longrightarrow\Phi(\Omega\otimes A/(f_1\ldots f_n))}
\]
\deleted{est un isomorphisme pour \deleted{\unsure{tt}} \deleted{système de paramètres} \inserted{\unsure{suite|suites} $A$-régulière} $(f_1\ldots f_n)$ de $\mathcal O_{X,x}$ (où $\Phi$ est $\Phi_R$ ou $\Phi_{R'}$)}

\deleted{b)}
\[
\deleted{\deleted{\unsure{\operatorname{Hom}(R|\operatorname{Hom}}}\Phi_{R'}(\Omega)\xrightarrow{\sim}\Phi_R(\Omega)}
\]

\deleted{Sous ces conditions, $u$ est un \uline{isom.}}

% archive sheet number: 12
